\chapter{Reusing causal software across languages without changing the question}
\label{exec17:s9}

The commercial comparisons showed that a result can change meaning at an
interface. Here the question is what it takes to preserve a computation across
languages. The argument proceeds from exact branches, shared fits, resampling
and graph identity to numerical composition and conditional path laws. It then
examines retries, persistent failure and selection correction, because the
law of returned results can change even when every completed calculation is
accurate. These derivations determine what a reuse or porting decision must
demonstrate.

\section{What must a port preserve?}

An original implementation, a language binding and a new implementation of the
same published method are different assets. For a bank, their comparison begins
with a decision: which result will inform a rate change, eligibility rule or
servicing policy? The artifact that must survive may be a graph class, a vector
of treatment contrasts, a joint uncertainty law or an explicit refusal. A port
that reproduces a point estimate while changing its population or resampling
unit has not reproduced that decision problem.

The immediate proposal is to reuse original-language implementations where
their scientific profile and dependencies fit. A native port needs a separate
case: a demonstrated interface, deployment or resource constraint and an
equivalence argument for the specific computation being replaced. This chapter
develops those obligations from inspected R, Python, Java and C++ sources. It
does not claim that the proposed workers have been built or that similar
algorithm names establish parity.

Let $F_v(x,u)$ be version $v$ of a source computation, with scientific input
$x$ and explicit random input $u$. A candidate $G_w$ in another language accepts
an encoding $E(x)$; its result is decoded by $D$. An exact coupled reference
would require
\begin{equation}
 D\{G_w(E(x),T(u))\}=F_v(x,u)
 \label{exec17:s9-pathwise}
\end{equation}
for every admitted $x,u$, where $T$ maps the random inputs. The equality must
include result status, target identity and uncertainty, not just numeric fields.
It is often stronger than a useful implementation needs. A numerical profile
can instead specify a justified error bound for a named output, while a
stochastic profile can require equality or an appropriate comparison of joint
output laws. These alternatives need to be selected explicitly; agreement on
one seed is not evidence of any universal equality.

The four languages expose different engineering choices, but no language
determines scientific validity. R pcalg offers source-visible FCI/RFCI paths
and conversion routines. Python CausalML exposes target columns, shared control
fits and row-bootstrap implementation. Java Tetrad stores graph endpoints
attached to nodes and implements distinct search branches. Native aGrUM
contains FCI/PAG code and is also accessible through bindings. The proposal
should preserve these originals as comparison references, then decide whether
to call, bind, port or reject the selected profile. Source inspection is the
evidence available here; it remains distinct from an upstream runtime trial
\citep{exec13_latentr_primary,exec14_causalml_primary,
exec14_java_latent_sources,exec14_native_agrum}.

\section{A branch comparison is part of the method}

The inspected pcalg Possible-D-Sep code has a concrete branch: after a supplied
conditional-independence test, it removes an edge when $p\geq\alpha$. A
missing test value is first replaced by the numeric value of the
\texttt{NAdelete} option. The same path has search-size and conditioning-order
limits. Translating the outer FCI name while changing any of these branches
can change the graph, even with identical input data
\citep{exec13_latentr_primary}.

Suppose a proposed numerical replacement returns $\widetilde p$ with
$|\widetilde p-p|\leq\epsilon$. If $p\leq\alpha-\delta$ or
$p\geq\alpha+\delta$ for $\delta>\epsilon$, the threshold decision agrees.
At an exact tie, replacing $\geq$ by $>$ disagrees even when $\epsilon=0$.
At $\alpha=1/20$, the values .0499, .05 and .0501 give delete decisions
$(0,1,1)$ under the inspected rule and $(0,0,1)$ under a strict comparison.
This finite reference concerns branch semantics; it does not establish that a
particular dataset produces those test values.

The margin argument extends to an adaptive algorithm only with care. At each
step, suppose the source and replacement have the same current graph, visit
the same candidate test, and their numeric error is bounded by $\epsilon_j$.
If the source comparison is separated from its threshold by more than
$\epsilon_j$ and ties/order are handled identically, the next branch agrees.
Induction then preserves the entire trace for that input. Without those
conditions, a first disagreement can change the next conditioning set, so a
bound on the original sequence of test values no longer controls the replacement
sequence. A maximum error in a final matrix is not a substitute for this
trace-level argument.

The appropriate public/generated fixture should therefore retain the visited
node pair, conditioning set, test status, effective threshold and branch
decision. Include exact ties, near ties and missing values. With
\texttt{NAdelete=true}, the inspected replacement of a missing value is 1,
which triggers deletion at the stated threshold; with false it is 0, which
does not. That is a source-defined policy for an unavailable test, not evidence
of independence or dependence. A bridge should retain the unavailable-test
diagnostic even when it reproduces the source branch. A new policy that refuses
the graph is permissible as a separately declared wrapper behavior, but it
must not be reported as exact source parity.

Search caps need similar treatment. The pcalg path can skip a Possible-D-Sep
search because the candidate set exceeds its configured maximum. Increasing
that maximum in a C++ port changes the algorithm's effective search, although
it may make sense for a new profile. Java Tetrad and aGrUM have their own
inspected traversal, filtering and orientation paths; their source details
cannot be collapsed into ``FCI with the same alpha.'' Candidate generation,
conditioning subsets, order, selection assumptions and stopping rules all
belong in a matched-method claim. The proposed first reuse route is the
original implementation with explicit controls. A port follows only after
the exact selected profile has a meaningful comparison oracle
\citep{exec13_latentr_primary,exec14_java_latent_sources,exec14_native_agrum}.

\section{Preserving the joint random computation}

CausalML's inspected T-learner code sorts non-control treatment labels, fits
one control model, and exposes that model as a shared reference for each
treatment group. This is useful computational reuse. It also means that
uncertainty in the control fit can be shared across treatment contrasts.
A port that independently refits the control for each action has changed
the joint computation, even if each individual contrast has the same marginal
distribution under an illustrative law \citep{exec14_causalml_primary}.

To isolate the issue, let $U_A,U_B,U_C$ be independent, equally probable
$-1$ and $+1$ errors in three fitted response means. Define two treatment
contrasts using one control:
\begin{equation}
 D_A=U_A-U_C,\qquad D_B=U_B-U_C.
 \label{exec17:s9-shared-control}
\end{equation}
Each has mean zero and variance two; their covariance is one. The difference
$D_A-D_B=U_A-U_B$ has variance two because the control error cancels. A
replacement using independent $U_{C,A},U_{C,B}$ instead has contrasts
$D'_A=U_A-U_{C,A}$ and $D'_B=U_B-U_{C,B}$. Each marginal is unchanged,
but their covariance is zero and the difference variance is four.
Enumerating all eight shared-control tuples and all sixteen split-control
tuples verifies
\begin{equation}
 \operatorname{Cov}\{(D_A,D_B)^\top\}=
 \begin{pmatrix}2&1\\1&2\end{pmatrix},\qquad
 \operatorname{Cov}\{(D'_A,D'_B)^\top\}=
 \begin{pmatrix}2&0\\0&2\end{pmatrix}.
 \label{exec17:s9-joint-covariance}
\end{equation}
The example is an authored random law, not a claimed variance formula for
CausalML or an observed vendor comparison.

A bank comparing two offers needs the joint contrast or a justified conservative
alternative, not merely two separately correct marginal intervals. Reusing a
fitted control object can preserve shared randomness; serializing it as two
independent fit instructions does not. Conversely, an implementation that
intentionally refits each contrast might be appropriate under its declared
method, but it is a different comparator. The source-derived reuse decision
should name object sharing, fit data, random state and target rows together.

Equal integer seeds in R, NumPy, Java and C++ do not by themselves establish
equal random draws. Even within one implementation, changing iteration order
can allocate the same random stream to different units or model components.
A proposed coupled comparison should supply explicit resample indices,
fold assignments and permitted tie-breaking draws whenever the original API
supports that control. Randomness inside a learner needs a separate identity
and rule. When exact coupling is unavailable, compare the required joint law
under a justified statistical design; do not infer failure from ordinary
Monte Carlo differences or infer success from two equal averages.

The distinction is particularly important for a proposed original-language
worker. Passing a model object by reference can retain sharing within one
process, while a worker boundary may require an immutable fitted artifact or
an explicit fit identifier. The common request can refer to that identity
without prescribing a generic Python pickle or Java object stream. The
scientific obligation is to preserve which fit is shared. Transport and
serialization are implementation choices to qualify after that obligation
has been stated.

\section{Rows, clusters and the law of a bootstrap}

The inspected CausalML bootstrap draws row indices with replacement. One
helper accepts a supplied seed for its random-index generator; another accepts
an RNG or uses NumPy's global state. Both slice features, treatment and outcome
with the same index vector. That alignment is essential, but it does not turn
row resampling into household, panel or calendar resampling
\citep{exec14_causalml_primary}.

Consider two observed account clusters, each containing two identical rows.
The first cluster has outcomes $(0,0)$, the second $(2,2)$. The empirical mean
is one. A four-row bootstrap samples from $(0,0,2,2)$ independently four
times. Its conditional mean remains one and its mean variance is $1/4$.
A two-cluster bootstrap instead draws two whole clusters from the two available
clusters and retains both rows in each selected cluster. Its mean also has
conditional expectation one, but variance $1/2$. The first result follows
from averaging four variance-one empirical draws; the second from averaging
two such independent cluster means.

The finite distributions are especially transparent. Four-row resampling has
$4^4=256$ equally likely index tuples; cluster resampling has four equally
likely cluster tuples and mean values $(0,1,1,2)$. Agreement in the returned
mean conceals a factor-of-two variance change. Neither conditional bootstrap
distribution is automatically a valid sampling approximation for bank data.
The example shows that silently changing the resampling unit during a port
is a change in the uncertainty procedure, even before asking whether that
procedure is scientifically justified.

For longitudinal servicing, a language boundary should therefore carry the
declared sampling unit and resampling identity, not infer them from the
number of rows in a data frame. If the original API supports only row
resampling, the first adapter can explicitly support independent-row targets
and refuse the requested dependent uncertainty. A new whole-account or
calendar resampling method needs its own assumptions and validation. It
should be presented as a method extension rather than silently enabled as
``the same bootstrap in R.'' Changing a seed cannot repair an incorrect
sampling unit.

Failed resamples are part of that contract. A bootstrap sample can omit a
treatment arm; deleting those replicates changes the conditional law from
the declared row bootstrap. Stratifying by arm also changes the law and may
be useful only as a declared procedure. The comparison needs all planned
replicates, valid-result indicators, diagnostics and aggregation rules.
A port that returns an interval from a selected successful subset may look
more reliable than a source that exposes failure, while answering a different
question. The proposed acquisition decision is to preserve the original
resampling path first and qualify any new dependent-data method separately.

\section{Endpoint types and stable identities}

A graph transfer has to preserve where each mark is attached. Write an edge
between stable node identifiers $i$ and $j$ as
\begin{equation}
  e_{ij}=(m_i,m_j),\qquad
  m_i,m_j\in\{\circ,\text{tail},\text{arrow}\}.
  \label{exec17:s9-edge-contract}
\end{equation}
The subscripts name the endpoint, not the array position. Thus a circle at $i$
and an arrow at $j$ remains a different object from an arrow at $i$ and a
circle at $j$. A complete graph object also contains the node universe, absent
edges, isolated nodes, edge multiplicity rules and any separate selection or
latent-couple annotations. Those items are not recoverable from a list of
directed arcs alone. An unresolved endpoint is a scientific result, not a
request for the receiving library to choose a convenient direction.

The inspected sources illustrate why a common label vocabulary is
insufficient. The selected pcalg export covers seven endpoint pairs and
therefore supplies no basis to assume that an arbitrary endpoint object round
trips. Tetrad endpoints are enum members attached to graph nodes; the selected
lookup uses node identity and returns null if the requested node is absent,
despite a comment that promises an exception. aGrUM's PAG code exposes a mark
at a destination node, while its generic mixed-graph conversion checks
arrowheads and can collapse circle/tail distinctions. A bidirected PAG edge
becoming two arcs is another information change. This is a warning about
specific inspected helpers, not a defect claim about every aGrUM API
\citep{exec13_latentr_primary,exec14_java_latent_sources,exec14_native_agrum}.

Define a versioned, partial translation $\phi_{v\to w}$ from the source
representation to the receiving representation. It is valid only on a stated
domain $\mathcal G_{v,w}$, and on that domain must satisfy
\begin{equation}
  \operatorname{decode}_w\!\left(\phi_{v\to w}(G)\right)=G
  \quad\text{up to a declared bijection of stable node identifiers.}
  \label{exec17:s9-round-trip}
\end{equation}
The equality includes marks at both endpoints, edge existence, isolated nodes,
and every annotation the requested operation needs. If the receiver has no
representation for a star, null, circle or selection annotation, the transfer
must return an explicit unsupported-result status. Mapping an unfamiliar
endpoint to a circle, dropping a latent-couple list, or fabricating a missing
edge is not a neutral fallback. A partial map is useful: it defines which
requests are safe to send and gives a precise refusal boundary.

The first comparison suite should generate all node-label permutations for
small labelled graphs, include isolated and reordered nodes, and exercise each
endpoint pair supported by both source and destination. It must distinguish
an absent edge from an edge whose mark is unknown. It should round-trip every
admitted graph, then assert refusal on each unsupported case. A graph that
matches only after sorting internal integer indices has passed a serialization
convenience test, not the label-permutation invariant. Stable identifiers must
be explicit and unique; display names alone may be duplicated or normalized
by a library. A source API keyed by object identity must first resolve each
object against the saved node map, before the graph is serialized.

This contract also prevents a false interpretation of package conversion
helpers. A helper may be correct for the subset it implements, while a caller
uses it outside that subset. The admission response therefore carries the
source graph class, receiver class, translated and unrepresented marks,
node-map digest, and any lost annotations. The downstream effect query can
then refuse to treat a partially represented graph as an identified set.
Graph conversion is not itself causal identification. Even an exact endpoint
round-trip preserves only the supplied graph object, not the assumptions that
made the graph scientifically credible.

\section{Observation clocks and row identity}

The row-bootstrap example has an analogue in the design matrix itself. The
inspected Tetrad lag utility creates rows by shifting positions in a data table;
its selected method accepts the dataset and lag count, not a subject ID or
physical timestamp. For a bank panel ordered as
\[
  (A,t_1),(A,t_2),(B,t_1),(B,t_2),
\]
a one-row lag assigns the preceding account's $t_2$ record to account B at
$t_1$. That transition is not part of either customer's history. The source
method can still be appropriate for one regularly sampled series, but the
port must not silently reuse it on stacked accounts. This boundary follows
from the inspected row construction; it is not a claim that every Tetrad time
series route lacks subject awareness \citep{exec15_javatemporal_tetrad}.

Before a worker receives a temporal table, the caller must state the
observation unit, subject key, timestamp or ordered event index, gap rule,
calendar/irregularity policy, and treatment of resets and overlapping windows.
The adapter must either construct within-subject sequences under that rule or
refuse to create lags. The model clock and inferential sampling unit are
separate fields. A customer is not automatically the sampling unit: uncertainty
may instead resample branches, campaigns, calendar blocks or randomized
clusters. Row counts cannot substitute for either concept.

For funding, an event ordered at origination cannot be joined to an outcome
that occurs before the funding decision merely because a later table contains
it. For servicing, a time-varying action requires its decision timestamp,
available history and censoring or competing-event status. These are caller
scientific inputs; a bridge can check their declared shape and order, but
cannot certify that a timestamp is truthful or an account ID correctly linked.
The response must retain the declared clock and preprocessing map so a
reviewer can tell whether a package's row position matches the intended bank
history.

\section{Which original should be reused?}

The evidence supports a narrow portfolio of original-language references, not
one universal engine. These are prospective route choices based on inspected
source, not successful local builds, matched runtime results or bank
validation.

For latent discovery in R, pcalg is the direct first reference for the selected
FCI/RFCI behavior because the pinned source exposes the search, missing-test,
threshold and endpoint-conversion paths. The proposed route calls a pinned R
worker with explicit controls and retains its diagnostics. The package declares
GPL (at least version 2), requires R 3.5 or later, links Rcpp/RcppArmadillo/BH
and imports additional graph packages. Dependency resolution and distribution
rights for the complete environment remain to be checked. A C++ rewrite is not
justified by a wish for one runtime; it first needs an actual deployment
constraint and must pass the trace and endpoint suite
\citep{exec13_pcalg_reuse_primary,exec13_latentr_primary}.

For heterogeneous data, rCausalMGM is a separate R candidate, not a drop-in
replacement for pcalg. Its inspected FCI-Stable documentation specifies a
mixed-data initial graph and orientation options, but also says that rows with
missing values are dropped; its selected profile and assumptions must travel
with the response. Its native implementation is wrapped for R through Rcpp
and declares GPL-3 plus Rcpp, survival, BH, Armadillo and RcppThread
requirements. Admit it only for an explicit complete-case target profile
until missingness handling and the exact mixed-data test are separately
qualified. It is not a standalone C++ library contract merely because some
algorithm code is in C++ \citep{exec14_native_rcmgm}.

MIIC is another distinct model family. The inspected package presents R and
C++ sources, a C++14 compiler requirement, Rcpp, and OpenMP flags delegated to
the R build. It declares GPL (at least version 2) and exposes information
theoretic, latent, and stationary temporal modes. Preserve MIIC as a comparative
candidate where its source assumptions and output semantics match the
question; do not translate its `latent' option into a generic PAG or report
its graph as an effect estimate. A direct C++ port is not established by the
inspected R wrapper. Scale and application claims are not measurements of the
target banking workload \citep{exec14_native_miic}.

For uplift-style effect estimation in Python, CausalML's pinned release is a
useful original reference for a selected binary T-learner profile. The root
license notice is Apache-2.0, but the source declares Python 3.11 or later and
eighteen direct requirements; transitive dependencies and deployment rights
have not been closed. Keep model fits and row bootstrap inside its Python
worker, carry shared-control fit identity explicitly, and treat learner
randomness separately from the resampling generator. This does not admit
dependent account-level uncertainty, and other estimators do not automatically
share this contract. Matching a CATE column without target rows and shared fit
does not preserve the bank contrast \citep{exec14_causalml_primary}.

For Java graph work, pinned Tetrad is a candidate where its graph/search profile
is needed and its GPL/dependency terms are accepted. The selected library POM
declares Java release 21 and a substantial dependency set. The same pinned
README describes Python and R entry routes as new and nonstandard, with a Java
25 requirement. This is a separate wrapper environment, not a contradiction
to resolve by guessing. Invoke one pinned artifact or one pinned wrapper stack,
record which was used, and do not mix them as interchangeable. Stable node IDs
and endpoint marks belong in the request before either path is admitted
\citep{exec13_tetrad_reuse_primary,exec14_java_latent_sources}.

For C++, aGrUM's native FCI/PAG path is a plausible comparator when deployment
calls for a compiled graph worker. Its inspected traversal, orientation and
endpoint methods identify a distinct implementation profile. The source
includes a dual LGPL-3-or-later/MIT statement for aGrUM, alongside warnings
that third-party code has separate copying terms. The review has not resolved
those dependencies, built the pinned library or qualified its PAG against
another implementation. Test the native profile on its own merits through a
narrow adapter; do not claim that the C++ function ports pcalg FCI because both
are called FCI \citep{exec14_native_agrum}.

These choices distinguish three actions often conflated in a software survey:
adopt an original implementation behind its native worker; bind existing
components while preserving their artifacts; or port selected logic only after
a measured need and bounded equivalence domain. Begin with original R and
Python workers for selected profiles, keep Tetrad and aGrUM as cross-check
candidates, and retain MIIC and rCausalMGM as distinct discovery families. No
package is rejected as scientifically useless; each is rejected as a universal
substitute where the inspected source supports a narrower claim. Keep native
reimplementation deferred until its benefit exceeds maintenance, validation,
rights and migration costs \citep{exec17_s9_route_decisions}.

\section{An admission contract for each worker}

The proposed boundary is a versioned request/response contract, not a generic
claim that packages can be hidden behind one function. The request names the
target population, action or graph question, outcome and units, information
cutoff, observation clock, row and sampling units, variable ontology,
missingness policy, identification assumptions, method controls, seed sources
and coupling, and requested uncertainty object. A source profile adds exact
package and dependency identities, representable graph/effect domain, timeout
and resource limits. Fields absent from a source API must be marked unsupported
or applied by a reviewed pre/postprocessor; they must not be silently ignored.

The response returns request identity, status, target identity, point or set
result, uncertainty type and sampling law, source/runtime identity, translated
node map or fit IDs, diagnostics, omitted rows/replicates, effective controls
and structured refusal. A successful process exit does not mean a scientific
result exists. An empty confidence set, unsupported graph mark, failed
resample, timeout and missing assignment field are distinct outcomes, not null,
zero or an empty list.

Conformance has five levels. First, schema checks establish request and unit
completeness. Second, exact fixtures establish representations, branch
conditions, row maps and unsupported cases. Third, numeric fixtures test
declared tolerances on admitted regular cases. Fourth, randomized outputs are
compared under shared-input coupling or a predeclared joint-distribution
comparison. Fifth, the pinned environment is installed, built, restarted and
faulted under recorded resource limits. A pass at one level does not imply a
pass at the next. The p-value tie and shared-control laws are gates for selected
levels, not production qualification.

For adaptive graph search, retain candidate edges, conditioning sets, CI-test
status, effective threshold and decisions. For effect workers, retain action
labels, target rows, fit reuse, resampling indices, failed replicates and
aggregation. For temporal workers, retain pre-lag row maps and subject/calendar
boundaries. Run the same cases through pinned source and bridge. If the native
package cannot expose an intermediate, output-level comparison may still be
possible, but missing diagnostics are a blind spot and cannot support a branch
parity claim.

Freeze acceptance criteria before comparisons. Exact discrete fields and
refusal statuses require exact agreement on the declared domain. Numeric fields
need a justified tolerance and a recorded consequence near decision
boundaries. Distributional criteria specify joint quantities and reference
law before draws. Operational criteria include reproducible installation,
assigned-request success, latency, memory, restart and failure accounting.
Unsupported requests and failed calls remain in the denominator. A retry that
retains only successful cases cannot inflate reported availability.

An evidence manifest pins source release, environment lock, compiler/JDK/R/
Python versions, build logs, transitive dependency identities, license texts,
source-to-output hashes and exact fixtures. Security and supply-chain review
are separate from semantic parity. A notice does not clear a transitive
license, and a process boundary alone does not settle redistribution rights.
Resolved environments and bank deployment terms are missing, so these
acceptance claims remain open.

\section{A banking consequence of shared computation}

Suppose retention must choose between offers A and B against one shared
control. Let $\widehat\theta_A$ and $\widehat\theta_B$ be incremental margins
and $k_A-k_B$ their operational cost difference. The decision uses
\[
  \widehat\Delta=\widehat\theta_A-\widehat\theta_B-(k_A-k_B).
\]
The shared control error can cancel in this difference; a port that fits
independent controls discards that covariance. Under the illustrative
Rademacher law, the variance is 2 under a shared fit and 4 under split fits.
More directly, the shared law puts probability $1/4$ on a positive contrast
difference; the split law puts $5/16$ there. These stipulated laws are not
estimates of CausalML. They show why matching marginal means may still change
a threshold or regret decision; exact example and enumeration are recorded
separately from upstream execution \citep{exec17_s9_bank_oracle}.

If a bridge preserves only marginal CATE vectors, it cannot reconstruct which
control fit was shared or whether joint covariance changed. It must return
marginals with a warning or refuse the paired choice; it must not imply
independent errors. First replay a public/generated request with a shared fit
identifier and exact target rows, then examine a predeclared dependent-bank
simulation. Port only if that route changes a decision-relevant measure and the
native API cannot satisfy the operating constraint.

This provides a falsifiable reason to integrate: preserve one specified bank
target and its uncertainty, not build a unified API for its own sake. Evidence
may favor the native worker, explicit refusal or a narrow port. It does not
assume one language or bridge will win \citep{exec17_s9_admission}.

\section{A joint-law boundary for port decisions}
\label{exec17:s9-law-boundary}

The shared-control example gives a decision-relevant reason to retain joint
randomness. Let $Z_A$ and $Z_B$ denote the two offer scores after subtracting
the same control score, and compare a bridge that independently redraws its
control for each offer. Under the finite Rademacher construction, exhaustive
enumeration gives
\begin{equation}
 \Pr(Z_A-Z_B=d)=
 \begin{cases}1/4&d=-2,\\1/2&d=0,\\1/4&d=2,
 \end{cases}\qquad\text{(shared control)},
 \label{exec17:s9-law-shared}
\end{equation}
and
\begin{equation}
 \Pr(\widetilde Z_A-\widetilde Z_B=d)=
 \begin{cases}1/16&d\in\{-4,4\},\\1/4&d\in\{-2,2\},\\3/8&d=0,
 \end{cases}\qquad\text{(split controls)}.
 \label{exec17:s9-law-split}
\end{equation}
The total variation distance of these two laws is $1/8$, and the probabilities
of a strictly positive difference are $1/4$ and $5/16$. A bridge can therefore
change the probability of selecting A even though both individual offer
contrasts have the same marginal variance. These are exact authored laws,
not a variance estimate or an execution of CausalML~\citep{exec17_s9_law}.

The corresponding numerical rule is also local to a declared decision. If a
source score $s$ and replacement score $\widetilde s$ satisfy
$|\widetilde s-s|\leq\epsilon$, then a threshold decision agrees whenever
$|s-c|>\epsilon$. At equality the bound cannot certify the branch. For a
finite action set with source values $v_a$, replacement errors bounded by
$\epsilon_a$, the source maximizer $a^*$ remains selected when
\begin{equation}
 v_{a^*}-\max_{a\ne a^*}v_a
   >\epsilon_{a^*}+\max_{a\ne a^*}\epsilon_a.
 \label{exec17:s9-action-gap}
\end{equation}
Without that gap, the replacement may choose a different action; its regret
is nevertheless bounded by $2\epsilon_{\max}$ under a uniform error bound.
This makes a porting claim testable at the level of the bank decision rather
than at an arbitrary number of decimal places. A result close to the boundary
must carry the interval or return an indeterminate-selection status.

For a language bridge, the finite law and the deterministic gap rule form two
different admission tests. The first asks whether shared inputs, fit objects
and resampling units preserve a joint output law. The second asks whether a
known numerical discrepancy can alter a selected action. Passing one does not
repair the other. The protocol should report both, together with the source
population, output status and refusal boundary. A bridge that cannot expose
shared-fit identity can still return marginal estimates, but it cannot claim
the paired offer decision is equivalent. Likewise, a graph endpoint map that
rounds a near-threshold p-value cannot call an adaptive search path identical
without the trace and margin evidence above.

The extension narrows the proposed integration decision. Native workers remain
the default when they satisfy the declared profile. A port earns a trial only
when its operational constraint is demonstrated, its admitted domain is
bounded, its exact and joint-law fixtures pass, and its decision-boundary
errors are disclosed. Otherwise the scientifically correct output is a
refusal or a native-worker recommendation. No software count or common API
can replace this evidence~\citep{exec17_s9_admission,exec17_s9_route_decisions}.

\section{Composing a numerical bridge across an adaptive calculation}
\label{exec17:s9-composition}

The decision to port one routine becomes more demanding when its output is
an input to the next routine. Consider a servicing analysis that transforms
observations, fits an outcome model, evaluates an offer contrast, and uses that
contrast to choose a subsequent diagnostic. Comparing only the final mean
discards the mechanism by which a small arithmetic difference can change the
diagnostic selected. Conversely, requiring identical machine bytes at every
stage can reject a useful implementation without showing any difference in
the decision. A proposal for a bridge needs a compositional statement between
these extremes. The following constructions are project derivations; their
exact finite references do not execute an upstream causal package
\citep{exec17_s9_composition}.

First fix the scientific input and a particular sequence of branches. Let the
source states satisfy $z_i=f_i(z_{i-1})$, for $i=1,\ldots,m$, and let the decoded
bridge states be $\widetilde z_i=\widetilde f_i(\widetilde z_{i-1})$. A state
includes the quantities needed by later stages, not merely the displayed
estimate. All norms and units below must be specified: a bound on a probability
cannot be inserted as if it were a bound on a monetary contrast. Suppose both
states remain in a declared domain $D_{i-1}$, the source map has Lipschitz
constant $L_i$ there, and the bridge has a uniform local discrepancy
$\|\widetilde f_i(z)-f_i(z)\|\leq\epsilon_i$ for $z\in D_{i-1}$.
If $\|\widetilde z_0-z_0\|\leq e_0$, the triangle inequality gives
\begin{equation}
 e_i=L_i e_{i-1}+\epsilon_i,\qquad
 \|\widetilde z_m-z_m\|\leq
 e_0\prod_{i=1}^{m}L_i+
 \sum_{i=1}^{m}\epsilon_i\prod_{j=i+1}^{m}L_j.
 \label{exec17:s9-composed-error}
\end{equation}
The empty product is one. The proof proceeds by comparing
$\widetilde f_i(\widetilde z_{i-1})$ first with
$f_i(\widetilde z_{i-1})$, then with $f_i(z_{i-1})$. Induction proves the
recursion and its expanded form. The domain condition is substantive. A source
trajectory inside a well-conditioned region does not show that its perturbed
trajectory stays there. The implementation needs a containing domain or an
explicit failure when the next enclosure leaves it.

For a transparent arithmetic example, take $z_0=0.12$,
$f_1(z)=2z+1$ and $f_2(z)=z/10$. The source result is $0.124$. Let input
discrepancy be at most $0.001$, stage-one error at most $0.002$, and
stage-two error at most $0.0001$. The composed bound is $0.0005$, attained
when these three errors have the positive sign. Simply adding the three
error allowances gives a different and less informative quantity because it
ignores amplification and contraction. These constants are stipulated for
the example. They are not certificates for a learner, a matrix inverse or
a floating-point library.

An adaptive search adds a second obligation. Suppose stage $i$ branches on
$g_i(z_{i-1})\geq c_i$, where the source comparison function has Lipschitz
constant $M_i$ and the bridge's direct comparison error is at most $\eta_i$.
The common branch is certified only if
\begin{equation}
 |g_i(z_{i-1})-c_i|>M_i e_{i-1}+\eta_i.
 \label{exec17:s9-composed-branch}
\end{equation}
This condition must be applied inductively on the actual common history.
After an uncertified branch, the two runs can ask different questions, so the
next same-function comparison may not exist. One cannot continue adding local
errors as if the branch sequence had remained fixed. The appropriate output is
a set of possible paths, a return to the original engine, or a declared lack
of equivalence. Equality at the boundary also requires preserving the source's
strict or non-strict comparison convention.

This distinction matters for discovery because a conditional-independence
decision determines which later conditioning sets are inspected. It also
matters for model selection: an apparently negligible change in a diagnostic
can select a different learner with a different uncertainty construction.
For these cases a useful porting experiment retains the intermediate state
enclosures and branch margins. Comparing final graphs without these traces
does not explain why agreement occurred, while a differing graph does not by
itself identify an arithmetic defect. The source may itself be close to a
decision boundary. The reuse choice should therefore depend on a declared
profile of stable paths, not a promise that all data admit the same numerical
certificate.

\section{Conditional laws, hidden state and adaptive calls}
\label{exec17:s9-adaptive-kernels}

Randomized software requires a statement about the law of a sequence of calls.
Let $H_i=(X_1,\ldots,X_i)$ be the decoded history after call $i$, including
statuses, fitted-object identities, branch choices and any persistent random
state needed by the next call. On finite state spaces write
$K_i(x\mid h)$ and $\widetilde K_i(x\mid h)$ for the source and bridge
conditional distributions of the next result given the same history $h$.
If the initial histories have the same law and every reachable conditional
kernel agrees, multiplication of conditional probabilities gives the same
joint path law. Equality of one-call marginals is insufficient because it
does not determine these conditional probabilities.

For example, a source draws one fair bit $B$ and returns $(B,B)$ in two
related calls. A bridge starts a fresh worker for each call and returns two
independent fair bits. Each call has the same marginal distribution. The
source returns agreeing bits with probability one; the bridge does so with
probability one half. Direct enumeration gives joint total variation distance
$1/2$. The missing state is the shared bit. The analogy in a bank analysis is
a fitted control model or common bootstrap draw that must be reused across
offer contrasts. Assigning a seed to each request is not enough if the
definition of a request has silently split a shared computation.

For approximate kernels define
$\delta_i=\sup_h \frac12\sum_x|K_i(x\mid h)-\widetilde K_i(x\mid h)|$,
where the supremum covers all common histories reachable while comparing the
two implementations. On a given common history, the distributions can be
paired so that they return the same value with probability
$\sum_x\min\{K_i(x\mid h),\widetilde K_i(x\mid h)\}$, at least
$1-\delta_i$. Put their common masses together, then distribute the residual
masses to complete this pairing. Repeating this construction until the first
disagreement yields
\begin{equation}
 d_{\rm TV}\{\mathcal L(H_m),\mathcal L(\widetilde H_m)\}
 \leq 1-\prod_{i=1}^{m}(1-\delta_i)
 \leq\sum_{i=1}^{m}\delta_i.
 \label{exec17:s9-path-tv}
\end{equation}
The first inequality follows because every event has the same indicator when
the coupled paths agree. The second follows by expanding the probability of
at least one disagreement. This is a sufficient bound, often conservative;
it does not assert that a measured discrepancy in a few requests bounds the
supremum over histories. If the initial laws differ, an additional initial
coupling factor is needed.

Adaptive stopping fits the same argument when the horizon is bounded in
advance. Add a stopped state that repeats itself through the remaining calls,
and include the stopping decision in the kernel. A cap on conditioning-set
size or on iterations then has an explicit role. With an unbounded horizon,
one additionally needs a summable error allocation or a controlled tail for
the number of calls. Stopping after apparently favorable numerical results
cannot be justified by applying a fixed-horizon bound only to completed paths.

The path bound and the arithmetic bound answer different questions. Tiny
shifts of continuous numerical outputs can have disjoint point masses and
total variation one even when their metric distance is tiny. For continuous
states, one can instead couple numerical discrepancies in a specified metric
and use the margin condition to certify a common discrete branch. To bound
the probability of a branch error, one must supply the probability mass within
the relevant margin and the probability that the arithmetic certificate fails.
A density or margin bound is an extra assumption, not a consequence of
observing a small error on one input. Thus the proposal should retain both the
numerical enclosure and the conditional law of the discrete choices, using
each only on the domain where its assumptions have been established.

\section{When retries select a different statistical result}
\label{exec17:s9-retry-selection}

A language bridge can preserve every successful arithmetic call and still
change the reported statistical law through its failure policy. Fix one
scientific request. Let $Y$ denote the result that an ideal computation with
random input would produce, with law $P$. Let $S$ indicate whether that attempt
completes and returns a usable result, and write
$q(y)=\Pr(S=1\mid Y=y)$. This representation does not assume a failed
attempt actually calculated or logged its latent $Y$. It describes the joint
law required to reason about selection. Assume $0<\bar q=\int q(y)P(dy)$.
Among successful attempts the reported law is
\begin{equation}
 P_{\rm success}(dy)=\frac{q(y)P(dy)}{\bar q},\qquad
 E(Y\mid S=1)-E(Y)=\frac{\operatorname{Cov}(Y,q(Y))}{\bar q},
 \label{exec17:s9-success-tilt}
\end{equation}
when the displayed expectations exist. Both identities follow by integrating
the joint success mass $q(y)P(dy)$. A completed-only comparison estimates the
first quantity. It estimates the intended source mean only under further
conditions such as constant completion probability, or zero covariance for
that particular mean. Zero covariance does not establish equality of whole
distributions.

Now allow at most $K$ attempts, redrawing both computational and operational
randomness independently each time under the same law. Return the first
successful result and otherwise return the explicit symbol $\bot$. Summing
over the position of the first success gives
\begin{align}
 \Pr(R\in dy)&=q(y)P(dy)\sum_{j=0}^{K-1}(1-\bar q)^j,
 \nonumber\\
 \Pr(R=\bot)&=(1-\bar q)^K,
 &E(N)&=\sum_{j=0}^{K-1}(1-\bar q)^j,
 \label{exec17:s9-redraw-law}
\end{align}
where $N$ counts attempted calls, including the final failed call when all
attempts fail. Conditional on a return, the result still has law
$P_{\rm success}$ for every $K$. More redraws improve completion but do not
remove outcome-dependent selection. Taking an unlimited number of attempts
can make completion approach one while retaining the same biased returned
law. If the same outage, memory pressure or data difficulty persists between
attempts, the independent-attempt premise itself fails.

Preserving computational randomness produces another procedure. Draw $Y$ once
and retry the operational execution of that same result. Suppose failures are
conditionally independent given $Y$, with the same $q(Y)$ on each attempt.
The chance that this result is eventually returned is
$q_K(y)=1-(1-q(y))^K$, and
\begin{align}
 \Pr(R\in dy)&=q_K(y)P(dy),\qquad
 \Pr(R=\bot)=\int(1-q(y))^K P(dy),
 \nonumber\\
 E(N)&=\sum_{j=0}^{K-1}\int(1-q(y))^jP(dy).
 \label{exec17:s9-frozen-law}
\end{align}
Completed results are weighted by $q_K$, not by $q$. If $q(y)>0$ almost
surely, this selection vanishes as $K$ grows without limit, although expected
cost can be unbounded. If $q(y)=0$ on a positive-mass region, those results
never appear. Retaining a seed cannot repair a deterministic lack of memory,
an unsupported data type or a nonterminating algorithm. Even conditional
independence can be inappropriate when retrying retains the same defective
process state.

Take a stipulated score error $Y\in\{-1,1\}$ with equal probabilities,
$q(-1)=1/4$, and $q(1)=1$. This is a selection example, not an empirical
claim about causal software. With two attempts, exhaustive enumeration of
signs and four equally likely operational states gives
\begin{center}
\begin{tabular}{lccc}
\hline
Procedure & Completion & $\Pr(Y=1\mid\text{return})$ & Returned mean\\
\hline
One attempt & $5/8$ & $4/5$ & $3/5$\\
Two independent redraws & $55/64$ & $4/5$ & $3/5$\\
Two retained-result attempts & $23/32$ & $16/23$ & $9/23$\\
\hline
\end{tabular}
\end{center}
Both two-attempt procedures have expected call count $11/8$, yet their
returned distributions and completion probabilities differ. A proposal that
specifies only the retry count and average cost leaves the statistical
procedure undetermined. If one score unit represents USD~100 of net offer
value, the ideal mean is zero while completed-only means are USD~60 and about
USD~39.13. These figures demonstrate a possible reporting distortion; they are
not estimates of the value of any banking intervention
\citep{exec17_s9_composition}.

\section{Persistent difficulty and honest all-request comparisons}
\label{exec17:s9-persistent-failures}

The independent redraw law can be optimistic even when each attempt has the
same marginal success rate. Let $\Theta$ denote a persistent request or
machine regime. Conditional on $\Theta=\theta$, assume attempts are
independent with success probability $q_\theta$. For a randomly assigned
request the no-return probability is
\begin{equation}
 \Pr(R=\bot)=E\{(1-q_\Theta)^K\}
 \geq (1-Eq_\Theta)^K,\qquad K\geq2.
 \label{exec17:s9-persistent-retry}
\end{equation}
The inequality follows because $x\mapsto(1-x)^K$ is convex on $[0,1]$.
It compares a mixture of conditionally independent attempt sequences with
an artificial sequence that redraws its regime on every attempt. In the
extreme finite example, half the regimes always complete and half never do.
Every attempt has marginal completion probability one half, yet no number
of within-regime retries reduces the no-return mass below one half. The
independent-attempt formula would instead give $2^{-K}$.

For a banking pipeline, difficult histories may create persistent large
conditioning sets, rare treatment arms or repeated optimizer failures.
Changing the language runtime may change this difficulty-selection relation.
Therefore the all-assigned comparison should retain the fixed request set,
its completion indicators, intended retry budget, actual attempts, and the
meaning of randomness retained between attempts. A useful analysis separates
variation between requests from repeated operational attempts within one
request. A thousand retries on one easy request are not a thousand
independent tests of the admitted request population.

There is a direct nonidentification issue if failed results cannot be
reconstructed. Fix one latent target result per assigned request, as in the
retained-result procedure. Suppose a bounded target $Y\in[a,b]$, completion fraction
$r$, and completed mean $\mu_s$ are observed under one declared procedure.
Writing $\mu_f$ for the unobserved failed-result mean gives
\begin{equation}
 E(Y)=r\mu_s+(1-r)\mu_f
 \in[r\mu_s+(1-r)a,\ r\mu_s+(1-r)b].
 \label{exec17:s9-failure-bounds}
\end{equation}
Without additional restrictions, any value in this interval can be attained
by assigning an admissible mean to failed requests. This is an identification
bound, not a confidence interval: sampling error in $r$ and $\mu_s$ needs
its own treatment. If the target is unbounded, finite bounds do not follow
from completion accounting alone. A timeout cannot honestly be coded as
zero, and a successful-subset mean cannot be labeled all-request performance.
For the retained-result example, the interval is $[0,9/16]$, containing
the ideal mean zero at its lower endpoint. Applying this decomposition to
independent redraws while calling the target the first draw's result would
be invalid: a returned later draw is a different latent quantity.

These mean bounds concern a latent analysis result. An operating decision
also needs the actual fallback action when analysis fails. Let $V$ denote
the underlying payoff state, distinct from the randomized analysis result
$Y$. If $U(A,V)$ is
the value of the chosen action and $C$ the total computation cost, the
relevant service quantity is
\begin{equation}
 E\{S_K U(A_{\rm returned},V)+(1-S_K)U(A_{\rm fallback},V)-C\},
 \label{exec17:s9-service-value}
\end{equation}
with $S_K$ indicating completion under the specified retry scheme. Comparing
only returned scores omits fallback consequences and cost; treating every
failure as the worst possible value can also distort the comparison if an
effective existing policy remains available. A public-data prototype can
stipulate these values transparently and vary them in sensitivity analysis.
It cannot infer a bank's actual fallback value, outage frequency or lost
business from the stipulation.

\section{Correcting selection, and the price of the correction}
\label{exec17:s9-selection-correction}

When the completion law is known and positive, one can recover some intended
functionals without claiming that the bridge has reproduced the source law.
For retained-result retries let $q_K(y)>0$ and let $h$ be the functional of
interest. Count every originally assigned request and use the contribution
$W=S_K h(Y)/q_K(Y)$, assigning zero on failure without needing its unobserved
$Y$. Conditional expectation gives
\begin{equation}
 E(W)=E\{h(Y)\},\qquad
 \operatorname{Var}(W)=E\{h(Y)^2/q_K(Y)\}-[E\{h(Y)\}]^2,
 \label{exec17:s9-inverse-selection}
\end{equation}
when the second moment is finite. The denominator of the empirical average
is the number originally assigned, not the number completed. This identity
requires a callable or independently justified $q_K$ for returned results.
It cannot be estimated from completed records alone without further
information about failures.

In the two-attempt retained-result example, $q_K(-1)=7/16$ and $q_K(1)=1$.
The inverse-weighted mean has expectation zero, but its per-request variance
is $23/14$, compared with one for an ideal observed score. Thus selection
correction restores the mean under its premises while increasing the
information cost. If the negative-score completion probability is mistakenly
entered as $1/2$ instead of $7/16$, the corrected expectation becomes $1/16$.
This residual bias is calculable even though every successful numeric score
is exact. It is not a floating-point error.

More generally, replacing $q_K$ by $\widehat q_K$ creates bias
$E[h(Y)\{q_K(Y)/\widehat q_K(Y)-1\}]$. If $|h|\leq H$,
$\widehat q_K\geq\varepsilon>0$ and
$|q_K-\widehat q_K|\leq d$ uniformly, its absolute value is at most
$Hd/\varepsilon$. A small absolute error in a rare completion probability
can therefore create a large error in the correction. Finite variance also
needs control of $h^2/q_K$, stronger than pointwise positivity. A practical
alternative is to narrow the admitted profile, make failed results observable
through a reference worker, or report bounds. Clipping inverse weights changes
the target unless its induced bias is separately bounded.

First-success redraws require another normalization. Given a returned result
with law proportional to $q(y)P(dy)$,
\begin{equation}
 E_P h(Y)=
 \frac{E_{\rm success}\{h(Y)/q(Y)\}}
      {E_{\rm success}\{1/q(Y)\}}.
 \label{exec17:s9-redraw-weighting}
\end{equation}
The numerator equals $E_P h/\bar q$ and the denominator equals $1/\bar q$,
which proves the identity. A finite-sample ratio of averages is generally
not unbiased, even with the correct $q$. Applying the retained-result
formula to independent-redraw outputs confuses two distinct sampling laws.
The retry semantics must consequently be part of the statistical procedure,
not an interchangeable configuration beneath it.

The resulting acquisition rule is specific. An original worker is attractive
when it permits retaining scientific state and exposing failures without
altering the procedure. A new bridge is attractive only when its proposed
benefit survives a comparison of numerical composition, conditional joint
laws, failure selection and fallback consequences. If completion is unrelated
to the target and branches remain certified, the obligations can simplify.
If these premises are unknown, a public-data experiment can test a bounded
profile and estimate diagnostics, while the proposal records the remaining
assumptions. No finite example here establishes universal port equivalence,
qualified dependent-data uncertainty, or a deployed banking service.
The next chapter uses these distinctions to compare build, buy and integration
routes on the same assigned workload and full lifecycle cost. Preserving the
procedure is an admission condition for that comparison, not a conclusion
that a port will be the least expensive choice.
